\section{Spiritual}

The spiritual dimension is not a metaphor here, it is a direction. The bulk is four-dimensional, and three of those dimensions show up on the cusp as the space we move in, while the fourth, the radial depth into the interior, is hidden from physics entirely. Almost all of the bulk lies off the cusp, unseen, so there is genuine room for an interior life and for structure the physical surface never shows. The spiritual is that depth, and the inward turn is movement along it, not an escape from the world but a descent into the same world's hidden thickness.

The depth wears several faces at once, and they coincide. It is the literal radial axis, the direction off the flat surface that no instrument on the surface can point along. It is the ladder of integration, since the deeper patterns are exactly the more unified ones, so that to go inward is to climb toward greater wholeness, the great chain of being and the radial axis turning out to be one line read two ways. It is the hidden web, the vast off-cusp structure that is real and addressed and yet never reaches the surface, the room where unseen things can live without any magic. And it is, on the one premise, the felt interior itself, the inner volume of experience that introspection turns to attend.

For there are two ways to move in the bulk. One is sideways, along the surface, parallel to the cusp, the ordinary travel of thought from one idea to its neighbor. The other is inward, along the radial axis, off the surface and into the depth, and that is the spiritual move. It is the same motion whether you call it climbing or diving, for integration is depth, and the climber toward the summit and the diver toward the abyss arrive at the same place measured in opposite directions. Going inward takes the self where the deepest and most integrated patterns sit, and where two things far apart on the surface are suddenly near.

The practices are specific motions along that axis. Introspection is attention turned inward, off the surface into the depth. Concentration is holding a deep region steady against the churn that would disperse it. Fluency in the dock's own directions is a kind of literacy in the moves of the bulk, and the meditation on turning, the felt fact that one full turn inverts a thing and only a second turn brings it home, is the double cover known from the inside. Equanimity is resting at the still tone, the zero, the center that is neither pole. And devotion, or surrender, is the alignment of one's own arrow with the one current of the whole, the ceasing to fight upstream. The last two are the most spiritual, because the conserved truth of the whole field is peace, and the deepest practice is to rest at that still center while the poles of pleasure and pain are held lightly and seen from their midpoint.

A self is also held apart from the world by its boundary, the screen that makes it a separate subject at all, and when that screen thins, by stillness or shock or long practice, the self feels its edge dissolve and its inside and the world's flow into one. This is the union the mystics report, placed exactly as a boundary gone slack and the self merging upward into a larger whole. Love runs the same way, two selves not merely near but integrating into one larger self, or phase-locking into resonance, or entering the same region of feeling-space so that each feels with the other rather than only about them.

Good and evil enter as the two poles of the one tone, pleasure and pain raised to the scale of a life, with peace the still center from which both are seen. All of this rests on the one premise, that the tone is felt, and it is offered as interpretation built on that premise, not as anything the bench can test.
